\section{The original closed-subscheme stalk and native ambient order}

\begin{definition}[The true original principal stalk quotient]
\label{def:fc-closed-immersion-principal-stalk-quotient}
\leanok
\lean{EnvelopeProperty.ClosedImmersionStalkQuotients.principalStalkQuotientEquiv, EnvelopeProperty.ClosedImmersionStalkQuotients.principalStalkQuotientEquiv_mk}
For the same actual closed immersion $g:D\hookrightarrow X$, an actual affine
open $U\subset X$, and an actual original equation $c$ generating the original
kernel ideal on $U$, the quotient of $\mathcal O_{X,g(y)}$ by the original germ
of $c$ is identified with the same original stalk $\mathcal O_{D,y}$.
The identification sends every quotient class of $a$ to the literal original
stalk map $g_y^\#(a)$.
\end{definition}

\begin{theorem}[True stalk lengths and native original orders]
\label{thm:fc-closed-immersion-stalk-order-length}
\leanok
\lean{EnvelopeProperty.ClosedImmersionStalkQuotients.stalkMap_ker_eq_map, EnvelopeProperty.ClosedImmersionStalkQuotients.principalStalkQuotient_length, EnvelopeProperty.ClosedImmersionStalkQuotients.ord_eq_closedStalk_length, EnvelopeProperty.ClosedImmersionStalkQuotients.residueFieldMap_isIso}
\uses{def:fc-closed-immersion-principal-stalk-quotient, thm:fc-scheme-order-local-length}
The kernel of the same original stalk map is the localization of the original
affine kernel ideal. Thus the length of the principal ambient stalk quotient
is exactly the length of the same closed-subscheme stalk over itself.
If $X$ is integral and locally Noetherian, $g(y)$ has codimension one, and the
original equation is nonzero, its native original function-field order at
$g(y)$ equals this true closed-subscheme local length. Moreover the literal
original residue-field map $\kappa(g(y))\to\kappa(y)$ is an isomorphism.
\end{theorem}
\begin{proof}
On the actual inverse-image affine open, closed immersion gives surjectivity
of the original section map. Both original germ maps are genuine prime
localizations, and the actual prime ideals correspond under the original
section map. Surjectivity maps the source prime-complement monoid onto the
target one, so the actual localization-kernel theorem identifies the whole
stalk kernel with the localization of the original affine kernel. For the
actual principal ideal, the first isomorphism theorem gives the displayed
whole quotient isomorphism, with its actual original element formula.
Restricting scalars along the same surjective stalk map preserves module
length, giving the exact length equality. The already proved native
order--principal-stalk-length formula gives the original function-field
order. Surjectivity of the same stalk map and of the actual residue quotient
makes the original residue-field map surjective; its source is a field, so
it is an isomorphism.
\end{proof}
